\section{A Punishment That Does Not End}

Every earthly affliction is endured within a life that passes. We may not know when relief will come, but time can bring a change, and death itself puts an end to temporal suffering. Hell has no last day. The punishment does not exhaust itself, finish a purification, or reach a point at which the condemned can begin another earthly trial. Christ calls it everlasting.

St.~Augustine attends closely to the final sentence of Matthew 25. Everlasting punishment and life everlasting stand within the same declaration, describing the two outcomes of judgment. He argues that the endless life of the saints and the endless punishment of the condemned must both be received from that sentence (\emph{City of God} XXI.23). The ground of his certainty is the truthfulness of the divine word. He will not set a conjectured eventual release above what Christ has said.

The chapter first considers the devil and his angels, to whom the same fire is appointed. Augustine then observes that the sentence joins the punishment of the wicked to theirs. His concern is the consistency of faith before a single Lord and Judge. The word which secures our hope of an everlasting heaven also teaches the permanence of hell.

Augustine also urges obedience while opportunity remains. The permanence of the revealed sentence gives the hearer a reason to attend to God's commands now. The argument has a practical consequence for someone who continues to dispute a duty he already recognizes. Further delay does not resolve the difficulty of obeying; it adds another refusal to those already made. Recollection can expose that repetition. The person needs to ask for grace to do what is required, even when he has not ceased to feel the attraction which first led him away.

The second death names this final ruin. In the thirteenth book, Augustine distinguishes bodily death, in which the soul leaves the body, from the death of the soul deprived of its good life with God. In final damnation, soul and body are reunited and the person continues to suffer. The soul animates the body, yet the whole person lacks blessed life with God. The condition is called death because of its penal misery, not because the creature ceases to exist (XIII.2).

This distinction changes the measure of our fear. Bodily death is a real enemy and a severe trial. Yet to avoid it by grave sin would place temporal survival above the life of grace. A person may lose earthly life in fidelity to Christ and gain everlasting life. Another may preserve health, security, and reputation while making choices that lead toward the second death.

The martyrdom of St.~Polycarp gives this ordering of fear a living form. The Church of Smyrna recounts his confession before the magistrate. When pressed to reproach Christ, he remembers the Lord he has served and refuses to blaspheme his King and Savior. Threats of beasts follow, then the threat of fire. Polycarp answers by distinguishing a fire that burns for a little time from the fire of coming judgment and eternal punishment (\emph{Martyrdom of Polycarp} 9--11).

Polycarp remembers Christ's faithful goodness when the magistrate demands blasphemy. His refusal answers a particular demand to betray his King and Savior. The threats do not change the character of that act, though they make obedience more painful. Smyrna's narrative then describes his confidence and joy (12). His answer about eternal fire explains why a threat of burning cannot compel his consent. He can lose bodily life while remaining faithful to the Lord who gives everlasting life.

The martyr's decision also clarifies choices in which the threatened loss is smaller. A profitable deception may seem difficult to abandon because others depend upon the income and the arrangement is longstanding. Those circumstances call for an honest way of meeting real obligations; they cannot make continued fraud right. The loss required by amendment belongs to this passing life. Polycarp's confession helps us judge that cost in relation to the good of remaining faithful.

St.~Alphonsus dwells on precisely this difference in his consideration on the eternity of hell. Even a limited pain becomes harder to bear as it continues, yet hell goes beyond every duration we can count. His argument rests on the scriptural teaching of everlasting punishment, then measures the fleeting pleasure of sin against that unending loss (\emph{Preparation for Death}, XXVII.1). Any number of years would still have an end; the revealed punishment has none.

The brevity of sinful pleasure often becomes evident even within this life. Anticipation magnifies it; possession soon passes. What remains may be guilt, damage to another, and the need to construct a further excuse. Eternity exposes the poverty of consenting to a grave evil for so little. It asks us to judge the promise of temptation by the end to which it leads, before its immediate attraction has taken command of the will.

St.~Cyprian brings finality and present mercy together in his address to Demetrianus. He describes enduring punishment of soul and body, with sorrow that can no longer bring fruitful repentance. Then he turns to the person he is addressing and urges him to seek salvation while life remains (\emph{Ad Demetrianum} 24--25). The opportunity will not remain forever; it has not yet ended. Neither age nor the number of past sins should prevent a living person from returning to God.

Cyprian addresses hostility toward Christians, yet offers his persecutors the way to rejoice with Christians eternally. He grounds this offered life in Christ's Cross, blood, and reconciliation of humanity to the Father. Their hostility has injured the people with whom he invites them to share everlasting joy. His appeal asks them to cease that hostility and receive the reconciliation Christ gives.

Cyprian's invitation addresses someone whose sins are already numerous and whose remaining life may be short. Neither fact prevents a present return to God. The years spent cannot be recovered, but their misuse can be repented of; the hostility practiced can cease. This is why the permanent sentence must be considered while we are still living. A sinner can now receive the mercy which Cyprian offers his persecutors. If shame over the past keeps him from returning, he is refusing the very reconciliation by which Christ heals that past.

